% The Sato-Tate law, seen live: a short Beamer talk with Sage code, for Sagebrush's LaTeX editor.
% Press Run to compute everything in this browser, then Present; the last slide but one has a live @interact.
% MIT license.
\documentclass{beamer}
\usetheme{Madrid}
\usepackage{sagetex}

\title[The Sato--Tate law]{How are the $a_p$ of an elliptic curve distributed?}
\subtitle{The Sato--Tate law, computed live}
\author{Sagebrush}
\date{November 2026}

\begin{document}

\begin{sagesilent}
E = EllipticCurve('37a1')
B = 10000

def normalized(C, bound):
    # x_p = a_p / (2 sqrt p) for the primes 3 <= p < bound where C has good reduction, sorted
    N = C.conductor()
    return sorted(float(C.ap(p)) / (2 * sqrt(float(p))) for p in prime_range(3, bound) if N % p)

def st_cdf():
    # the Sato-Tate distribution function on [-1, 1]
    return plot(1/2 + (x*sqrt(1 - x^2) + arcsin(x))/pi, (x, -1, 1), color='red', thickness=2)

def ecdf(xs):
    # the fraction of the xs at or below each of them, the points joined by lines
    return line([(t, (i + 1)/len(xs)) for i, t in enumerate(xs)], color='blue', figsize=[6, 3.4])

xs = normalized(E, B)
small = [p for p in prime_range(2, 30)]
counts = table([['$p$'] + small, [r'$\#E(\mathbb{F}_p)$'] + [p + 1 - E.ap(p) for p in small],
                ['$a_p$'] + [E.ap(p) for p in small]])
moments = table([(2*k, n(sum(t^(2*k) for t in xs)/len(xs), digits=3), catalan_number(k)/4^k) for k in [1, 2, 3]],
                header_row=['$2k$', r'mean of $x_p^{2k}$', '$C_k/4^k$'])
\end{sagesilent}

\begin{frame}
  \titlepage
\end{frame}

\begin{frame}{An elliptic curve}
  \[ E:\ y^2 + y = x^3 - x \]
  \begin{itemize}
    \item<2-> Conductor $N = \sage{E.conductor()}$, discriminant $\Delta = \sage{E.discriminant()}$.
    \item<3-> $j$-invariant $j = \sage{E.j_invariant()}$, not an integer, so $E$ has no complex multiplication.
    \item<4-> Rank $\sage{E.rank()}$: $E(\mathbb{Q})$ is infinite.
  \end{itemize}
\end{frame}

\begin{frame}{Counting points modulo $p$}
  For a prime $p \nmid N$, reduce $E$ modulo $p$ and count:
  \[ \#E(\mathbb{F}_p) = p + 1 - a_p. \]
  \pause
  \begin{center}
    \sage{counts}
  \end{center}
\end{frame}

\begin{frame}{Hasse's bound}
  Hasse (1936): \quad $|a_p| \le 2\sqrt{p}$, \quad so \quad $x_p = \dfrac{a_p}{2\sqrt{p}} \in [-1, 1]$.
  \pause
  \begin{center}
    \sageplot[width=0.85\textwidth]{list_plot([(p, E.ap(p)/(2*sqrt(p))) for p in prime_range(3, 2000) if p != 37], size=4, figsize=[8, 3.2], ymin=-1.05, ymax=1.05)}
  \end{center}
  \pause
  The points look random. Is there a law?
\end{frame}

\begin{frame}{The Sato--Tate law}
  \begin{block}{Theorem (Sato--Tate, proved 2008--2011)}
    If $E/\mathbb{Q}$ has no complex multiplication, the $x_p$ are equidistributed in $[-1, 1]$ with density
    \[ \mu(x) = \frac{2}{\pi}\sqrt{1 - x^2}. \]
  \end{block}
  \pause
  Writing $a_p = 2\sqrt{p}\cos\theta_p$, the angles $\theta_p \in [0, \pi]$ follow $\frac{2}{\pi}\sin^2\theta\,d\theta$.
  \pause
  \medskip

  Proved by Clozel, Harris, Shepherd-Barron and Taylor when $j$ is not an integer, and by Barnet-Lamb,
  Geraghty, Harris and Taylor for every curve without complex multiplication.
\end{frame}

\begin{frame}{Evidence: the moments}
  The even moments of $\mu$ are $\displaystyle\int_{-1}^{1} x^{2k}\,\mu(x)\,dx = \frac{C_k}{4^k}$, with $C_k$ the Catalan numbers.
  \pause
  \medskip

  For the $\sage{len(xs)}$ primes $p < \sage{B}$ with $p \nmid 2N$:
  \begin{center}
    \sage{moments}
  \end{center}
\end{frame}

\begin{frame}{Evidence: the distribution function}
  The fraction of the $x_p$ at most $t$, at each $x_p$ and joined by lines (blue), against the Sato--Tate
  distribution function (red)
  \[ F(t) = \frac12 + \frac{t\sqrt{1 - t^2} + \arcsin t}{\pi}. \]
  \begin{center}
    \sageplot[width=0.6\textwidth]{ecdf(xs) + st_cdf()}
  \end{center}
\end{frame}

\begin{frame}[fragile]{Complex multiplication is different}
\begin{sagesilent}
G = EllipticCurve('27a1')
ps = prime_range(5, B)
zeros = len([p for p in ps if G.ap(p) == 0])
\end{sagesilent}
  $G:\ y^2 + y = x^3 - 7$ has $j = \sage{G.j_invariant()}$ and complex multiplication by $\mathbb{Q}(\sqrt{-3})$.
  \pause
  \medskip

  For $p \equiv 2 \pmod 3$, $G$ is supersingular and $a_p = 0$: that is
  $\sage{zeros}$ of the $\sage{len(ps)}$ primes $5 \le p < \sage{B}$, about half.
  \pause
  \medskip

  So the $x_p$ of $G$ can't follow $\mu$: half of them sit at $0$.
\end{frame}

\begin{frame}[fragile]{Try it}
\begin{sagesilent}
@interact
def sato_tate(curve=selector(['11a1', '37a1', '389a1', '27a1', '32a1'], default='37a1'),
              bound=slider([1000, 2000, 5000, 10000], default=2000)):
    C = EllipticCurve(curve)
    ys = normalized(C, bound)
    print('{}: {} primes, mean of x_p^2 = {:.3f}'.format(curve, len(ys), sum(t^2 for t in ys)/len(ys)))
    show(ecdf(ys) + st_cdf())
\end{sagesilent}
  Pick a curve and a bound. The curves 27a1 and 32a1 have complex multiplication.
\end{frame}

\begin{frame}{References}
  \small
  \begin{itemize}
    \item H.~Hasse, Zur Theorie der abstrakten elliptischen Funktionenk\"orper I--III, \emph{J.~reine angew.\ Math.}\ 175 (1936).
    \item L.~Clozel, M.~Harris, R.~Taylor, Automorphy for some $\ell$-adic lifts of automorphic mod $\ell$ Galois representations, \emph{Publ.\ Math.\ IH\'ES} 108 (2008).
    \item R.~Taylor, Automorphy for some $\ell$-adic lifts of automorphic mod $\ell$ Galois representations~II, \emph{Publ.\ Math.\ IH\'ES} 108 (2008).
    \item M.~Harris, N.~Shepherd-Barron, R.~Taylor, A family of Calabi--Yau varieties and potential automorphy, \emph{Ann.\ of Math.}\ 171 (2010).
    \item T.~Barnet-Lamb, D.~Geraghty, M.~Harris, R.~Taylor, A family of Calabi--Yau varieties and potential automorphy~II, \emph{Publ.\ RIMS} 47 (2011).
    \item A.~V.~Sutherland, Sato--Tate distributions, \emph{Contemp.\ Math.}\ 740 (2019).
  \end{itemize}
\end{frame}

\end{document}
